%%=============================================================================
%% Samenvatting
%%=============================================================================

% TODO: De "abstract" of samenvatting is een kernachtige (~ 1 blz. voor een
% thesis) synthese van het document.
%
% Een goede abstract biedt een kernachtig antwoord op volgende vragen:
%
% 1. Waarover gaat de bachelorproef?
% 2. Waarom heb je er over geschreven?
% 3. Hoe heb je het onderzoek uitgevoerd?
% 4. Wat waren de resultaten? Wat blijkt uit je onderzoek?
% 5. Wat betekenen je resultaten? Wat is de relevantie voor het werkveld?
%
% Daarom bestaat een abstract uit volgende componenten:
%
% - inleiding + kaderen thema
% - probleemstelling
% - (centrale) onderzoeksvraag
% - onderzoeksdoelstelling
% - methodologie
% - resultaten (beperk tot de belangrijkste, relevant voor de onderzoeksvraag)
% - conclusies, aanbevelingen, beperkingen
%
% LET OP! Een samenvatting is GEEN voorwoord!

%%---------- Nederlandse samenvatting -----------------------------------------
%
% TODO: Als je je bachelorproef in het Engels schrijft, moet je eerst een
% Nederlandse samenvatting invoegen. Haal daarvoor onderstaande code uit
% commentaar.
% Wie zijn bachelorproef in het Nederlands schrijft, kan dit negeren, de inhoud
% wordt niet in het document ingevoegd.

\IfLanguageName{english}{%
\selectlanguage{dutch}
\chapter*{Samenvatting}
\lipsum[1-4]
\selectlanguage{english}
}{}

%%---------- Samenvatting -----------------------------------------------------
% De samenvatting in de hoofdtaal van het document

\chapter*{\IfLanguageName{dutch}{Samenvatting}{Abstract}}

In deze bachelorproef wordt onderzocht hoe apps voor stress- en burn-outpreventie het mentale welzijn van Vlaamse twintigers kunnen ondersteunen. Deze doelgroep wordt steeds sterker blootgesteld aan prestatiedruk, sociale verwachtingen en permanente digitale connectiviteit. Tegelijkertijd nam in België tussen 2018 en 2023 het aantal personen jonger dan 30 jaar dat langdurig arbeidsongeschikt is door burn-out met 107\% toe. Hoewel het aanbod aan digitale toepassingen voor mentaal welzijn sterk is gegroeid, blijven vragen bestaan over hun doelgroepgerichtheid, gebruiksvriendelijkheid en mogelijke cognitieve belasting.

\vspace{1em}

Het onderzoek startte met een literatuurstudie naar stress en burn-out bij twintigers, digitale mentale gezondheidsinterventies, preventiestrategieën, digitalisering en cognitieve belasting, en \acrshort{ui}/\acrshort{ux}-ontwerp. Op basis hiervan en van een interview werden functionele en niet-functionele requirements opgesteld volgens de MoSCoW-methode. Daarnaast werden ontwerpprincipes geformuleerd die beschrijven hoe digitale toepassingen mentale rust kunnen bevorderen en cognitieve belasting kunnen beperken. Deze requirements vormden de basis voor een vergelijkende analyse van bestaande mentale gezondheidsapps. Vijf geselecteerde toepassingen werden geëvalueerd aan de hand van de opgestelde ontwerpprincipes. Hieruit werden relevante functionaliteiten en zowel sterke als minder effectieve ontwerpkeuzes geïdentificeerd. Op basis van deze inzichten werd een cross-platform \acrshort{poc} ontwikkeld met onder meer gepersonaliseerde onboarding, een toolkit met oefeningen, dagboek- en reflectiefuncties, psycho-educatieve content en instelbare notificaties en voorkeuren.

\vspace{1em}

Uit het onderzoek blijkt dat een stressreducerende applicatie niet alleen moet focussen op de aangeboden interventies, maar ook op de manier waarop deze worden aangeboden. Belangrijke uitgangspunten zijn personalisatie, een overzichtelijke interface, beperkte informatie, digitale onderbrekingen, en het vermijden van mechanismen die extra prestatiedruk veroorzaken. Zo wordt aanbevolen om notificaties instelbaar te maken, informatie stapsgewijs aan te bieden, één primaire taak per scherm te hanteren en gamification zoals streaks en rankings te vermijden.

\vspace{1em}

De \acrshort{poc} werd ten slotte heuristisch geëvalueerd aan de hand van de eerder opgestelde requirements en ontwerpprincipes. De resultaten tonen aan dat deze principes bruikbaar zijn als praktisch kader voor het ontwerpen en beoordelen van stressreducerende digitale toepassingen. 

\vspace{1em}

Het onderzoek kent echter beperkingen, waaronder het ontbreken van een evaluatie met de uiteindelijke doelgroep van de app. Er kan daarom niet worden geconcludeerd dat de toepassing daadwerkelijk stress of burn-out vermindert. Verder onderzoek met Vlaamse twintigers is nodig om zowel de ontwerpprincipes als de \acrlong{poc} empirisch te valideren.
